\documentclass[12pt]{letter}
\usepackage[margin=1in]{geometry}
\usepackage{amsmath,amssymb}
\usepackage{hyperref}

\makeatletter
\newcommand{\paragraph}[1]{\medskip\noindent\textbf{#1}\enspace}
\makeatother

\signature{Andrew Craton\\
Research Director, Kootru Labs\\
Burlington, Wisconsin, USA\\
acraton@kootru.com\\
ORCID: 0009-0001-2269-8599}

\address{Andrew Craton\\
Kootru Labs\\
Burlington, Wisconsin, USA}

\begin{document}

\begin{letter}{Editorial Board\\
PRX Quantum\\
American Physical Society}

\opening{Dear Editors,}

I submit for your consideration the manuscript entitled
\textbf{``Provable measurement design for quantum state tomography
from Hamming-scheme spectral theory,''}
which establishes the first non-adaptive measurement
design for quantum state tomography with provable
reconstruction guarantees, derived from a new
correspondence with Hamming association schemes.
The manuscript is submitted for publication in PRX Quantum
(Article format, 16 pages of main text plus appendices).

\paragraph{Summary.}
Quantum state tomography in the underdetermined regime requires
choosing which Pauli operators to measure from an exponentially
large set.
Random sampling, the dominant approach, ignores operator weight
and wastes circuit budget on weight classes with little
reconstruction value.
This manuscript establishes a correspondence, new to both
combinatorics and quantum information, between Hamming
association schemes and quantum measurement design.

Krawtchouk polynomials on $H(n,2)$ decompose the Pauli operator
space into weight-class eigenspaces and fix a parameter-free
allocation: a single integer $K^* = 5$ determines the full
measurement set (137~operators at $n = 4$, compiling into
29~tensor-product circuits) with no optimization and no
target-state input.
In the reverse direction, the basin separation theorem
(Theorem~2) proves that \emph{weight saturation}, a
Hamming-scheme property, is sufficient for MLE basin separation:
every maximizer of the log-likelihood agrees on all informative
Pauli coordinates.
The spectral decomposition holds if and only if $q \le 3$
(a characterization new to both fields); the
weight-block-diagonal structure that drives basin separation
extends to all~$q$ (verified through $q = 7$).

On superconducting hardware, structured selection improves
reconstruction fidelity by $\Delta F = +0.33 \pm 0.07$ over
random sampling (IBM \texttt{ibm\_fez}, four independent runs;
corroborated at $+0.25$ on Rigetti Ankaa-3), growing to
$+0.53 \pm 0.06$ at eight~qubits under a compositional
architecture that uses $\le 116$~circuits for any linear chain length.
Across 23~test states (product, entangled, Haar-random, mixed),
$K^*$ ranks first among seven non-adaptive strategies including
D-optimal experimental design and derandomized classical shadows
at matched budget.
Applied to a Floquet discrete time crystal on IBM hardware,
$K^*$ recovers nearest-neighbor XY correlators consistent with
noiseless simulation ($r = 0.94$, 8/9 sign agreement) that are
structurally inaccessible to Z-basis-only protocols.

To our knowledge, this provides the first non-adaptive measurement
design for quantum state tomography with provable MLE
identifiability guarantees, replacing heuristic operator
selection with an allocation derived from algebraic first
principles.

\paragraph{Why PRX Quantum.}
This work addresses PRX Quantum's criterion for
\emph{exceptional connection} between quantum science and a
neighboring discipline:

\begin{itemize}
\item \emph{Forward (combinatorics $\to$ quantum).}
  Krawtchouk spectral analysis has established applications in
  quantum error correction and state transfer but has not
  previously been applied to measurement allocation.
  Here it produces the first non-variational operator-selection
  principle for underdetermined tomography that carries a provable
  reconstruction guarantee (Theorem~2).

\item \emph{Reverse (quantum $\to$ combinatorics).}
  The basin separation theorem gives weight saturation of
  Hamming-scheme subsets a reconstruction-theoretic rigidity not
  visible from classical distance properties, connecting
  association-scheme combinatorics to nonlinear estimation
  landscapes for the first time.
  This result appears new to coding theory independently of
  its quantum application.

\item \emph{Universal and operationally significant.}
  The allocation mechanism (eigenvalue masses determining
  per-weight budget; weight saturation determining reconstruction
  guarantee) extends to qudits of any local dimension~$q$ via the
  weight-block-diagonal structure, with full Krawtchouk
  diagonalization available for $q \le 3$.
  A compositional architecture delivers constant-circuit scaling
  ($\le 116$~circuits for any linear-chain length), and the
  advantage grows with system size ($+0.33$ at $n = 4$ to $+0.53$
  at $n = 8$) because random selection dilutes as $4^{-n}$ per
  weight class while the structured allocation is size-independent.
  As processors scale beyond 100~qubits, local tomography of
  subsystem density matrices becomes the primary characterization
  task; the deterministic, constant-cost $K^*$ protocol provides
  full local states, making it a natural measurement infrastructure
  for routine device diagnostics.

\item \emph{Timely application to many-body quantum phases.}
  Discrete time crystal experiments rely on Z-basis magnetization
  to certify period doubling (Mi et al., Nature \textbf{601}, 531,
  2022), but this observable cannot distinguish quantum from
  classical DTCs.
  $K^*$'s Hamming-weight allocation naturally resolves the conserved
  sectors of the XXZ interaction, providing access to off-diagonal
  correlators ($\ell_1$ coherence, partial-transpose negativity,
  XY terms) that encode the quantum character of time-crystalline
  order.
  A proof-of-principle run on IBM \texttt{ibm\_kingston} recovers
  nearest-neighbor XY correlators of a 4-qubit Floquet DTC that
  match noiseless predictions ($r = 0.94$).
\end{itemize}

Recent PRX Quantum publications have established similarly
productive cross-disciplinary bridges: chi-boundedness from graph
theory applied to quantum measurement (King et al., 2025),
algebraic graph theory applied to noise learnability
(Chen \& Flammia, 2026), and measurement frame theory applied
to shadow tomography (Innocenti et al., 2023).
The present work bridges a comparable gap between Hamming
association schemes and quantum tomography.
It differs from a one-directional import in that the basin
separation theorem exports a new result \emph{to} combinatorics,
and the correspondence is validated on multi-vendor hardware.

\paragraph{Durability beyond the present noise floor.}
As gate fidelities improve, the raw-fidelity advantage over random
sampling will narrow on generic states.
Three structural properties persist regardless of noise floor:
certifiable worst-case bounds via the basin separation theorem,
direct single-shot access to off-diagonal correlators, and a
compositional architecture whose circuit budget is constant per
patch.
In the fault-tolerant regime, single-patch reconstruction
($n = 4$--$9$) remains the core characterization primitive for
code-block verification, exactly the regime $K^*$ was
constructed for.

\paragraph{Validation and scope.}
The theory makes a testable prediction: GHZ states should yield
$F \approx 0.50$ due to weight-class mismatch.
Hardware confirms this exactly, and a weight-class fingerprint
classifier detects the condition automatically, recovering
$F = 0.887$ from the same data at zero additional circuit cost.
This predict-then-confirm cycle illustrates the explanatory
power of the spectral framework beyond raw fidelity improvement.

Experimental scope spans two superconducting vendors (IBM,
Rigetti), four state families, and qubit counts $n = 2$--$8$.
Simulation compares 23~states across seven non-adaptive
strategies at matched budget.
All raw data, reconstruction code, and the complete
137-operator measurement set are provided as Supplemental
Material, together with a verification pipeline spanning four
independent backends and Lean~4 formal proofs with
zero \texttt{sorry} obligations.

\paragraph{Suggested reviewers.}
Given the cross-disciplinary scope, I suggest one reviewer in algebraic
combinatorics or association schemes (Hamming schemes, Krawtchouk
polynomials, Delsarte theory) and one in quantum state tomography
or measurement design (classical shadows, optimal experimental
design, MLE reconstruction).

\paragraph{AI disclosure.}
Per APS policy: I formulated all conjectures, chose all high-level
proof strategies, made all substantive mathematical decisions, and
verified every proof with deterministic scripts (SymPy, NumPy);
no AI tools were used in the verification pipeline.
Claude (Anthropic, Opus~4.6) was used as a proof-formalization
tool: drafting LaTeX proofs from my specified conjectures and
strategies, and for copyediting manuscript prose.  Claude
occasionally selected specific algebraic techniques to execute
the strategies I specified.
All 11 formal statements were subsequently machine-verified in
Lean~4 (Mathlib, v4.29.0-rc8) with zero \texttt{sorry} obligations;
a CI pipeline enforces sorry-freeness on every commit
(Appendix, ``Lean\,4 formal verification status'').
Full details are in the ``Verification and AI disclosure'' section.

I confirm that no AI tools were used as co-authors, and I take
full responsibility for all content.  The manuscript has not been
submitted elsewhere.

\paragraph{Competing interests.}
The author declares no competing financial or non-financial interests.

\paragraph{Verification resources.}
Every numerical claim in the manuscript is reproduced
end-to-end by a 3-minute Jupyter notebook at\\
\url{https://github.com/kootru-repo/kstar-verify}\\
(\texttt{notebook/k\_star\_demo.ipynb}, runnable in-browser
via Binder with no local install); the tiered onramp in
\texttt{VERIFICATION.md} ranges from a 30-second manifest
peek through a 30-minute Lean rebuild from scratch.

\closing{Sincerely,}

\end{letter}
\end{document}
